\usepackage{tikz}
\usepackage{pdfpages}
\usepackage{xcolor}
\usepackage{float}
\usetikzlibrary{fit, calc, positioning, matrix, arrows.meta, shadows}

% Вариант 1: Создание команды \highlight
\newcommand{\highlight}[1]{\colorbox{yellow!50}{#1}}

% Создание шапки титульной страницы
\renewcommand{\titlepageheader}[2]
{
	\begin{wrapfigure}[7]{l}{0.14\linewidth}
		\vspace{3.4mm}
		\hspace{-8mm}
		\includegraphics[width=0.89\linewidth]{bmstu-logo}
	\end{wrapfigure}
	
	{
		\singlespacing \small
		Министерство науки и высшего образования Российской Федерации \\
		Федеральное государственное автономное образовательное учреждение \\
		высшего образования \\
		<<Московский государственный технический университет \\
		имени Н.~Э.~Баумана \\
		(национальный исследовательский университет)>> \\
		(МГТУ им. Н.~Э.~Баумана) \\
	}
	
	\vspace{-4.2mm}
	\vhrulefill{0.9mm} \\
	\vspace{-7mm}
	\vhrulefill{0.2mm} \\
	\vspace{2.8mm}
	
	{
		\small
		ФАКУЛЬТЕТ \longunderline{<<#1>>} \\
		\vspace{3.3mm}
		КАФЕДРА \longunderline{<<#2>>} \\
	}
}

% Установка заголовков РПЗ
\renewcommand{\titlepagenotetitle}[2]
{
	{
		\LARGE \bfseries
		РАСЧЕТНО-ПОЯСНИТЕЛЬНАЯ ЗАПИСКА \\
	}
	\vspace{5mm}
	{
		\Large \itshape
		#1 \\
		\vspace{5mm}
		НА ТЕМУ: \\
		\vspace{5mm}
		<<#2>> \\
	}
}

% Создание страницы определений
\renewenvironment{definitions}
{
	\chapter*{ОПРЕДЕЛЕНИЯ}
	\addcontentsline{toc}{chapter}{ОПРЕДЕЛЕНИЯ}
	
	В настоящей расчетно-пояснительной записке применяют следующие термины с соответствующими определениями.
	
	\begin{description}[leftmargin=0pt]
	}
	{
	\end{description}
}

% Создание страницы обозначений и сокращений
\renewenvironment{abbreviations}
{
	\chapter*{ОБОЗНАЧЕНИЯ И СОКРАЩЕНИЯ}
	\addcontentsline{toc}{chapter}{ОБОЗНАЧЕНИЯ И СОКРАЩЕНИЯ}
	
	В настоящей расчетно-пояснительной записке применяют следующие сокращения и обозначения.
	
	\begin{description}[leftmargin=0pt]
	}
	{
	\end{description}
}

\RequirePackage{enumitem}
\newcommand{\abbreviation}[2]
{
	\item \noindent #1 --- #2
}

\newcommand{\ssr}[1]
{
	\chapter{#1}
	\addcontentsline{toc}{chapter}{#1}  
}

\newcommand{\ssrno}[1]
{
	\chapter*{#1}
	\addcontentsline{toc}{chapter}{#1}  
}


% Переопределение формата даты обращения
\DeclareFieldFormat{urldate}
{%
	(дата~обращения:%
	\addspace
	\iffieldundef{urlday}{}{\ifnumless{\thefield{urlday}}{10}{0\thefield{urlday}}{\thefield{urlday}}%
		\adddot
		\iffieldundef{urlmonth}{}{\ifnumless{\thefield{urlmonth}}{10}{0\thefield{urlmonth}}{\thefield{urlmonth}}%
			\adddot
			\thefield{urlyear}}})
}

% Переопределение списков
\setlist[itemize]{label=---}

% Команды для различных типов слоев
\newcommand{\convlayer}[4][]{
	% Сверточный слой
	\foreach \i in {0,...,#4} {
		\node[layer, convcolor, right=2cm of #2, yshift=\i*(-6pt)+(#4*3pt), xshift=\i*(-6pt)+(#4*3pt), #1] (#3\i) {};
	}
}

\newcommand{\poollayer}[4][]{
	% Подвыборочный слой
	\foreach \i in {0,...,#4} {
		\node[smalllayer, poolcolor, right=2.5cm of #2, yshift=\i*(-6pt)+(#4*3pt), xshift=\i*(-6pt)+(#4*3pt), #1] (#3\i) {};
	}
}

\newcommand{\fclayer}[4][]{
	% Полносвязный слой
	\foreach \i in {0,...,#4} {
		\node[fclayer, fccolor, right=2.75cm of #2, yshift=\i*(-6pt)+(#4*3pt), xshift=\i*(-6pt)+(#4*3pt), #1] (#3\i) {};
	}
}

% Команда для капсульных слоев
\newcommand{\capsulelayer}[4][]{
	% Капсульный слой
	\foreach \i in {0,...,#4} {
		\node[layer, capsulecolor, right=2.5cm of #2, yshift=\i*(-6pt)+(#4*3pt), xshift=\i*(-6pt)+(#4*3pt), #1] (#3\i) {};
	}
}

% Команда для капсульных слоев
\newcommand{\attentionlayer}[4][]{
	% Капсульный слой
	\foreach \i in {0,...,#4} {
		\node[layer, attentioncolor, right=2.5cm of #2, yshift=\i*(-6pt)+(#4*3pt), xshift=\i*(-6pt)+(#4*3pt), #1] (#3\i) {};
	}
}

\newcolumntype{Y}{>{\centering\arraybackslash}X}

\makeatletter
\renewcommand*{\l@chapter}[2]{
	\ifnum \c@tocdepth>\m@ne
	\addpenalty{-\@highpenalty}
	\vskip 1em \@plus.2em
	\@dottedtocline{0}{0pt}{1.5em}{\bfseries #1}{\bfseries #2}
	\fi
}
\makeatother
